\documentclass[10pt,a4paper]{amsart}
\usepackage[margin=19mm]{geometry}
\usepackage{amsmath,amssymb,mathtools}
\usepackage[scr=rsfs,cal=euler]{mathalfa}
\usepackage{xfrac}
\usepackage[unicode,hidelinks,bookmarks=false]{hyperref}
\newcommand{\ie}{i.e.\ }
\newcommand{\cf}{cf.\ }
\newcommand{\resp}{respectively\ }
\newcommand{\Subsection}{\subsection}
\providecommand{\nobreakdash}{\nobreak\hbox{-}}
\setlength{\emergencystretch}{2em}
\multlinegap0pt
\usepackage{bm}
\usepackage{bbm}
\usepackage{dsfont}
\usepackage{xspace}
\usepackage{cancel}
\usepackage{mathtools}
\usepackage{amsthm}
\makeatletter
\@ifundefined{thm}{\newtheorem{thm}{Thm}}{}
\@ifundefined{lem}{\newtheorem{lem}[thm]{Lemma}}{}
\@ifundefined{prop}{\newtheorem{prop}[thm]{Proposition}}{}
\makeatother
\newcommand{\C}{\mathbb{C}}
\renewcommand{\Im}{\operatorname{Im}}
\renewcommand{\L}{\mathcal{L}}
\renewcommand{\Re}{\operatorname{Re}}
\renewcommand{\i}{\sqrt{-1}}
\title[The asymptotic behaviors of the colored Jones polynomials of]{The asymptotic behaviors of the colored Jones polynomials of the figure eight-knot, and an affine representation}
\author{Hitoshi Murakami}
\date{Source: arXiv:2307.07100v1 (CC BY 4.0)}
\begin{document}
\maketitle

\noindent\emph{Source and licence.} The asymptotic behaviors of the colored Jones polynomials of the figure eight-knot, and an affine representation. Version arXiv:2307.07100v1 (CC BY 4.0), distributed under the Creative Commons Attribution 4.0 International licence, as recorded in the arXiv licence metadata for this exact version and shown on the abstract page https://arxiv.org/abs/2307.07100v1. Full rights notice: licenses/arxiv-2307.07100-CC-BY-4.0.txt.\par\smallskip

\noindent\textit{Editorial scope.} Counted unit: the Prop whose source label is \texttt{prop:TN\_L2} in the article (source0.tex lines 1606--1617, source label \texttt{prop:TN\_L2}), delivered together with the article's own complete original proof of it (source0.tex lines 1618--1828, one \texttt{proof} environment of 211 lines). No mathematics is written, repaired or re-derived: statement and proof are the article's own text line for line. Required context retained uncounted: eq:TN\_1 (lines 526--530); eq:TN\_gamma (lines 565--569); eq:TN\_negative (lines 785--791); eq:TN\_positive (lines 795--801); lem:TN\_gamma\_Omega (lines 907--909); lem:L2\_L1\_gamma (lines 1179--1219); lem:L2\_L1 (lines 1501--1515); lem:TN\_L2 (lines 1592--1603). Every definition, display and label of those lines is retained unchanged. Omitted from this package and not redistributed: 1--525, 531--564, 570--784, 792--794, 802--906, 910--1178, 1220--1500, 1516--1591, 1604--1605, 1829--4648. Nothing inside a retained excerpt is omitted or altered: every retained body is delivered exactly as the article's lines are written, so no omission receipt of any kind accompanies this package. Citations: the retained excerpts carry no citation command at all, so no bibliography entry of the article is transcribed and no reference is redistributed. Reference adaptation: none was needed and none was made; every reference inside the delivered unit resolves inside it, so no unresolved reference or citation is produced. Printed slips kept literal: no printed word, symbol, hypothesis or number of the counted statement or of its proof was altered. Layout and class: the article's own class is \texttt{amsart} with packages including amssymb, mathrsfs, stmaryrd, verbatim, epic, eepic, hyperref, graphicx, tikz; the wrapper is the lane's pinned \texttt{amsart} editorial wrapper at 10pt on A4, which restates the theorem-like environments the retained text needs and adds the article's own macro definitions that the retained text needs (\texttt{\textbackslash C}, \texttt{\textbackslash Im}, \texttt{\textbackslash L}, \texttt{\textbackslash Re}, \texttt{\textbackslash i}), with extra wrapper packages (bm, bbm, dsfont, xspace, cancel, mathtools). The delivered numbering is the unit's own. Assets: no retained span contains a figure environment, a graphics-inclusion command, a PDF-page inclusion, a table or a TikZ picture; no image file is copied into this package, the package declares no asset dependency and no image file is redistributed with it. Licence: version arXiv:2307.07100v1 of this article is distributed under the Creative Commons Attribution 4.0 International licence, as recorded in the arXiv licence metadata for this exact version and shown on the abstract page https://arxiv.org/abs/2307.07100v1. A full rights notice is delivered at licenses/arxiv-2307.07100-CC-BY-4.0.txt. Characters: every retained excerpt is pure ASCII, so the delivered engine, XeLaTeX with its default Latin Modern text font, reports no missing character.
\begin{equation}\label{eq:TN_1}
  T_{N}(z)-T_{N}(z+1)
  =
  \L_1\left(\frac{N}{\gamma}z+\frac{1}{2}\right).
\end{equation}

\begin{equation}\label{eq:TN_gamma}
  T_{N}\left(z-\frac{\gamma}{2N}\right)-T_{N}\left(z+\frac{\gamma}{2N}\right)
  =
  \L_1(z).
\end{equation}

\begin{equation}\label{eq:TN_negative}
  T_{N}(z)
  :=
  T_{N}(z+1)
  +
  \L_1\left(\frac{N}{\gamma}z+\frac{1}{2}\right),
\end{equation}

\begin{equation}\label{eq:TN_positive}
  T_{N}(z)
  :=
  T_{N}(z-1)
  -
  \L_1\left(\frac{N}{\gamma}(z-1)+\frac{1}{2}\right),
\end{equation}

\begin{lem}\label{lem:TN_gamma_Omega}
The function $T_{N}(z)$ extended by using \eqref{eq:TN_1} satisfies \eqref{eq:TN_gamma} for $z\in\Omega$.
\end{lem}

\begin{lem}\label{lem:L2_L1_gamma}
Let $\nu$ and $M$ be positive real numbers with $0<\nu<1/2$.
Then, we have
\begin{equation*}
  \frac{N}{2\pi\i\gamma}\L_2\left(z-\frac{\gamma}{2N}\right)
  -
  \frac{N}{2\pi\i\gamma}\L_2\left(z+\frac{\gamma}{2N}\right)
  =
  \L_1(z)
  +
  O(N^{-2})
\end{equation*}
as $N\to\infty$ for $z$ in the region
\begin{equation}\label{eq:L2_L1_approx}
  \{z\in\C\mid-1+\nu\le\Re{z}\le2-\nu,|\Im{z}|\le M\}
  \setminus\left(\square^{-}_{\nu}\cup\square^{+}_{\nu}\right),
\end{equation}
where
\begin{align*}
  \square^{-}_{\nu}
  &:=
  \{z\in\C\mid-1+\nu\le\Re{z}\le\nu,|\Im{z}|\le\nu\},
  \\
  \square^{+}_{\nu}
  &:=
  \{z\in\C\mid1-\nu\le\Re{z}\le2-\nu,|\Im{z}|\le\nu\}.
\end{align*}
This means that there exists a constant $c>0$ that does not depend on $z$ such that
\begin{equation*}
  \left|
    \frac{N}{2\pi\i\gamma}\L_2\left(z-\frac{\gamma}{2N}\right)
    -
    \frac{N}{2\pi\i\gamma}\L_2\left(z+\frac{\gamma}{2N}\right)
    -
    \L_1(z)
  \right|
  <
  \frac{c}{N^2}
\end{equation*}
for sufficiently large $N$.
\end{lem}

\begin{lem}\label{lem:L2_L1}
There exist positive real numbers $c$ and $\varepsilon$ such that
\begin{equation*}
  \left|
    \frac{N}{2\pi\i\gamma}\L_2(z)
    -
    \frac{N}{2\pi\i\gamma}\L_2(z+1)
    -
    \L_1\left(\frac{N}{\gamma}z+\frac{1}{2}\right)
  \right|
  <
  ce^{-\varepsilon N}
\end{equation*}
for any $z$ in the region $\C\setminus\bowtie_{\nu}$ if $N$ is sufficiently large.
\end{lem}

\begin{lem}\label{lem:TN_L2}
Let $\nu$ and $M$ be positive real numbers.
Then, there exists a constant $c>0$ such that
\begin{equation*}
  \left|
    T_{N}(z)-\frac{N}{2\pi\i\gamma}\L_2(z)
  \right|
  =
  \frac{c}{N}
\end{equation*}
for $z$ in the region $\{z\in\C\mid\nu\le\Re{z}\le1-\nu,|\Im{z}|\le M\}$ if $N$ is sufficiently large, where $c$ does not depend on $z$.
\end{lem}

\begin{prop}\label{prop:TN_L2}
Suppose that $\nu<1/4$.
We have
\begin{equation*}
  T_{N}(z)
  =
  \frac{N}{2\pi\i\gamma}\L_2(z)
  +
  O(N^{-1})
\end{equation*}
as $N\to\infty$ in the region $\Omega^{\ast}_{\nu}$.
\end{prop}

\begin{proof}
We need to prove the proposition for $z$ with $-1+\nu\le\Re{z}<\nu$ or $1-\nu<\Re{z}\le2-\nu$.
\par
If $z\in\Omega^{\ast}_{\nu}$ and $-1+\nu\le\Re{z}<-\nu$, we use \eqref{eq:TN_negative}.
We have
\begin{align*}
  &\left|
    T_{N}(z)-\frac{N}{2\pi\i\gamma}\L_2(z)
  \right|
  =
  \left|
    T_{N}(z+1)
    +
    \L_1\left(\frac{N}{\gamma}z+\frac{1}{2}\right)
    -
    \frac{N}{2\pi\i\gamma}\L_2(z)
  \right|
  \\
  \le&
  \left|
    T_{N}(z+1)-\frac{N}{2\pi\i\gamma}\L_2(z+1)
  \right|
  \\
  &+
  \left|
    \L_1\left(\frac{N}{\gamma}z+\frac{1}{2}\right)
    -
    \frac{N}{2\pi\i\gamma}\L_2(z)
    +
    \frac{N}{2\pi\i\gamma}\L_2(z+1)
  \right|
  =
  O(1/N),
\end{align*}
where we apply Lemmas~\ref{lem:TN_L2} and \ref{lem:L2_L1}, noting that we can apply Lemma~\ref{lem:L2_L1} because $z\notin\bowtie_{\nu}$.
\par
Similarly, if $z\in\Omega^{\ast}_{\nu}$ and $1+\nu<\Re{z}\le2-\nu$, using \eqref{eq:TN_positive}, we have
\begin{align*}
  &\left|
    T_{N}(z)-\frac{N}{2\pi\i\gamma}\L_2(z)
  \right|
  =
  \left|
    T_{N}(z-1)
    -
    \L_1\left(\frac{N}{\gamma}(z-1)+\frac{1}{2}\right)
    -
    \frac{N}{2\pi\i\gamma}\L_2(z)
  \right|
  \\
  \le&
  \left|
    T_{N}(z-1)-\frac{N}{2\pi\i\gamma}\L_2(z-1)
  \right|
  \\
  &+
  \left|
    -\L_1\left(\frac{N}{\gamma}(z-1)+\frac{1}{2}\right)
    -
    \frac{N}{2\pi\i\gamma}\L_2(z)
    +
    \frac{N}{2\pi\i\gamma}\L_2(z-1)
  \right|
  =
  O(1/N),
\end{align*}
noting that we can apply Lemma~\ref{lem:L2_L1} because $z-1\notin\bowtie_{\nu}$.
\par
If $z\in\Omega^{\ast}_{\nu}$ and $-\nu\le\Re{z}<\nu$, we put $m:=\left\lfloor\frac{2N\nu}{\Re\gamma}\right\rfloor+1$.
From Lemma~\ref{lem:TN_gamma_Omega} we have
\begin{equation*}
\begin{split}
  T_{N}(z)
  =&
  T_{N}\left(z+\frac{\gamma}{N}\right)
  +
  \L_1\left(z+\frac{\gamma}{2N}\right)
  \\
  =&
  T_{N}\left(z+\frac{2\gamma}{N}\right)
  +
  \L_1\left(z+\frac{3\gamma}{2N}\right)
  +
  \L_1\left(z+\frac{\gamma}{2N}\right)
  \\
  =&\cdots
  =
  T_{N}\left(z+\frac{m\gamma}{N}\right)
  +
  \sum_{j=1}^{m}\L_1\left(z+\frac{(2j-1)\gamma}{2N}\right).
\end{split}
\end{equation*}
Now since $m\le\frac{2N\nu}{\Re\gamma}+1<m+1$, we have $\nu<\Re\bigl(z+\frac{m\gamma}{N}\bigr)<3\nu+\frac{\Re\gamma}{N}<1-\nu$ if $N>\frac{\Re\gamma}{1-4\nu}$, and so we can apply Lemma~\ref{lem:TN_L2} to $z+\frac{m\gamma}{N}$.
We have
\begin{align*}
  &\left|
    T_{N}(z)-\frac{N}{2\pi\i\gamma}\L_2(z)
  \right|
  \\
  =&
  \left|
    T_{N}\left(z+\frac{m\gamma}{N}\right)
    -
    \frac{N}{2\pi\i\gamma}\L_2(z)
    +
    \sum_{j=1}^{m}\L_1\left(z+\frac{(2j-1)\gamma}{2N}\right)
  \right|
  \\
  \le&
  \left|
    T_{N}\left(z+\frac{m\gamma}{N}\right)
    -
    \frac{N}{2\pi\i\gamma}\L_2\left(z+\frac{m\gamma}{N}\right)
  \right|
  \\
  &+
  \left|
    \frac{N}{2\pi\i\gamma}\L_2\left(z+\frac{m\gamma}{N}\right)
    -
    \frac{N}{2\pi\i\gamma}\L_2(z)
    +
    \sum_{j=1}^{m}\L_1\left(z+\frac{(2j-1)\gamma}{2N}\right)
  \right|
  \\
  =&
  \left|
    \frac{N}{2\pi\i\gamma}\L_2\left(z+\frac{m\gamma}{N}\right)
    -
    \frac{N}{2\pi\i\gamma}\L_2(z)
    +
    \sum_{j=1}^{m}\L_1\left(z+\frac{(2j-1)\gamma}{2N}\right)
  \right|
  +O(1/N).
\end{align*}
Since we have
\begin{equation*}
  \L_2\left(z+\frac{m\gamma}{N}\right)-\L_2(z)
  =
  \sum_{j=1}^{m}
  \left(
    \L_2\left(z+\frac{j\gamma}{N}\right)
    -
    \L_2\left(z+\frac{(j-1)\gamma}{N}\right)
  \right),
\end{equation*}
we have
\begin{equation}\label{eq:sum_L2_L1}
\begin{split}
  &\left|
    \frac{N}{2\pi\i\gamma}\L_2\left(z+\frac{m\gamma}{N}\right)
    -
    \frac{N}{2\pi\i\gamma}\L_2(z)
    +
    \sum_{j=1}^{m}\L_1\left(z+\frac{(2j-1)\gamma}{2N}\right)
  \right|
  \\
  \le&
  \sum_{j=1}^{m}
  \left|
    \frac{N}{2\pi\i\gamma}\L_2\left(z+\frac{j\gamma}{N}\right)
    -
    \frac{N}{2\pi\i\gamma}\L_2\left(z+\frac{(j-1)\gamma}{N}\right)
  \right.
  \\
  &
  \phantom{\sum_{j=1}^{m}}\hfill
  \left.
    +
    \L_1\left(z+\frac{(2j-1)\gamma}{2N}\right)
  \right|.
\end{split}
\end{equation}
We use Lemma~\ref{lem:L2_L1_gamma} to conclude that each summand of the right-hand side of \eqref{eq:sum_L2_L1} is less than $\frac{c}{N^2}$ for $c>0$.
Note that $c$ is independent of $j$.
Since $m=\left\lfloor\frac{2N\nu}{\Re\gamma}\right\rfloor+1$, the right-hand side of \eqref{eq:sum_L2_L1} is less than
\begin{equation*}
  \frac{mc}{N^2}
  \le
  \left(\frac{2N\nu}{\Re\gamma}+1\right)
  \frac{c}{N^2}
  \le
  \frac{c'}{N}
\end{equation*}
if we put $c':=\left(\frac{2c\nu}{\Re\gamma}+1\right)$.
\par
If $1-\nu<\Re{z}\le1+\nu$, from Lemma~\ref{lem:TN_gamma_Omega} we have
\begin{equation*}
\begin{split}
  T_{N}(z)
  =&
  T_{N}\left(z-\frac{\gamma}{N}\right)
  -
  \L_1\left(z-\frac{\gamma}{2N}\right)
  \\
  =&
  T_{N}\left(z-\frac{2\gamma}{N}\right)
  -
  \L_1\left(z-\frac{3\gamma}{2N}\right)
  -
  \L_1\left(z-\frac{\gamma}{2N}\right)
  \\
  =&\cdots
  =
  T_{N}\left(z-\frac{m\gamma}{N}\right)
  -
  \sum_{j=1}^{m}\L_1\left(z-\frac{(2j-1)\gamma}{2N}\right),
\end{split}
\end{equation*}
where we put $m:=\left\lfloor\frac{2N\nu}{\Re\gamma}\right\rfloor+1$ as before.
Since $\nu<\Re\bigl(z-\frac{m\gamma}{N}\bigr)<1-\nu$ as before, we can prove the proposition similarly.
\end{proof}

\end{document}
